% ===== latex-recomposed:document =====
% flm: style file flmrecomposed.b84bd50e.sty
\begin{theorem}\flmLabels{\label{ref1}}
A theorem no proof refers back to.
\end{theorem}

\begin{proof}[{\flmProofTitle{*nolabel}}]
A back-reference whose text is no label: its title is shown as written.
\end{proof}

\begin{proof}[*ref2]
A back-reference naming nothing.
\end{proof}%%
% ===== latex-recomposed:style flmrecomposed.b84bd50e.sty =====
\NeedsTeXFormat{LaTeX2e}
\ProvidesPackage{flmrecomposed.b84bd50e}
% --- packages ---
% flm.latex.pkg.phfthm
\RequirePackage[thmset=empty]{phfthm}
% flm.latex.pkg.zref-clever
\RequirePackage{zref-clever}
% --- commands and environments the library defines for the document ---
% flm.latex.def.theorems
% Theorems. \noproofref, written after a statement's labels, says that no proof
% of the statement is coming; defined here as nothing where no package (phfthm)
% defined it first.
\ProvideDocumentCommand\noproofref{}{}
% The kinds of a class are set in the theorem style of the class's name. The
% shipped class definitionlike is set in amsthm's definition style (upright);
% declare a theorem style of that name before loading this file to change it.
\ifcsname th@definitionlike\endcsname\else\newtheoremstyle{definitionlike}{}{}{\normalfont}{}{\bfseries}{.}{5pt plus 1pt minus 1pt}{}\fi
% flm.latex.def.theorems.authored
% \flmProofTitle{<title>}: what a proof's heading sets for a title the document
% wrote, which is written as its argument inside the proof environment's
% optional argument. phfthm reads that argument as the thing the proof proves
% and words the heading itself (Proof of <title>), so by default the title
% alone. Redefine \flmProofTitle to set a proof's title differently.
\ProvideDocumentCommand\flmProofTitle{m}{#1}
% flm.latex.def.labels
% \flmLabels{<\label commands>}: the labels of one place in the document.
% In text it runs them as written. On a line of a display-math
% environment, where amsmath keeps the last label of a line and loses the
% ones before it, it runs the first and records every further one as an alias:
% a line per alias is written to the .aux file after the last page and read on
% the next run, making the alias refer to what the first label refers to
% (for \ref and for \zcref alike). Redefine \flmLabels to change what a
% group of labels does; \def\flmLabels#1{#1} restores plain LaTeX behavior,
% amsmath's "Multiple \label's" message included.
\ProvideDocumentCommand\flmLabels{m}{%
  \ifmmode\expandafter\flm@labels@math\else\expandafter\@firstofone\fi{#1}}
% Reading the group: the first \label as written, the rest as aliases of it.
\def\flm@labels@math#1{\def\flm@lbl@first{}\flm@lbl@scan#1\flm@lbl@stop}
\def\flm@lbl@scan{\@ifnextchar\label\flm@lbl@take\flm@lbl@rest}
\def\flm@lbl@take\label#1{%
  \ifx\flm@lbl@first\@empty
    \def\flm@lbl@first{#1}%
    \label{#1}%
  \else
    \expandafter\flm@recordalias\expandafter{\flm@lbl@first}{#1}%
  \fi
  \flm@lbl@scan}
\def\flm@lbl@rest#1\flm@lbl@stop{#1}
% Collecting the aliases. A name is recorded once however often the line is
% typeset (amsmath measures a line before setting it).
\def\flm@aliaslist{}
\def\flm@recordalias#1#2{%
  \ifcsname flm@alias@#2\endcsname\else
    \expandafter\gdef\csname flm@alias@#2\endcsname{}%
    \g@addto@macro\flm@aliaslist{\flm@alias@entry{#2}{#1}}%
  \fi}
% Writing them, from the hook that runs after the last page has been shipped
% out: every \newlabel of that page has reached the .aux file by then, so the
% alias lines come after them. \AtEndDocument would run too early.
\def\flm@alias@write#1#2{\immediate\write\@auxout{\string\flm@aliaslabel{#1}{#2}}}
\AddToHook{enddocument/afterlastpage}{\let\flm@alias@entry\flm@alias@write\flm@aliaslist}
% Reading them back: the alias refers to what the original refers to. The
% assignments are global because the .aux file is read inside a group.
\def\flm@aliaslabel#1#2{%
  \ifcsname r@#2\endcsname
    \global\expandafter\let\csname r@#1\expandafter\endcsname\csname r@#2\endcsname
  \fi
  \ifcsname Z@R@#2\endcsname
    \global\expandafter\let\csname Z@R@#1\expandafter\endcsname\csname Z@R@#2\endcsname
  \fi}
% --- definitions written from the configuration of FLM features ---
% th@theoremlike (the theorems feature's, written from its configuration): declared unless \th@theoremlike is already defined, so a definition made earlier under that name is kept
\ifcsname th@theoremlike\endcsname\else\newtheoremstyle{theoremlike}{}{}{\itshape}{}{\bfseries}{.}{5pt plus 1pt minus 1pt}{}\fi
% theorem (the theorems feature's, written from its configuration): declared unless \theorem is already defined, so a definition made earlier under that name is kept
\ifcsname theorem\endcsname\else\phfMakeTheorem[thmstyle=theoremlike,counter=phfthmcounter]{theorem}{Theorem}\zcRefTypeSetup{theorem}{Name-sg={Theorem},name-sg={theorem},Name-pl={Theorems},name-pl={theorems},Name-sg-ab={Thm.},name-sg-ab={Thm.},Name-pl-ab={Thms.},name-pl-ab={Thms.}}\AddToHook{env/theorem/begin}{\zcsetup{countertype={theorem=theorem}}}\fi
% proof (the theorems feature's, written from its configuration): declared unless \proof is already defined, so a definition made earlier under that name is kept
\ifcsname proof\endcsname\else\phfMakeProofEnv[defaultproofname={Proof}]{proof}\fi
\endinput
